\chapter{Conclusions}
\label{chapter:conclusions}

Finally, \Cref{chapter:conclusions} summarizes the entire body of work presented in the thesis. This chapter will provide a summary of the key contributions and findings from each research thrust, reiterating the advancements made in applying deep learning to REIMS data for seafood authentication. It will discuss the overall significance of the work in the context of food science, analytical chemistry, and industrial applications. The limitations of the current research will be acknowledged, followed by a comprehensive outline of promising future work directions, suggesting avenues for further dataset expansion, model refinement, real-world deployment, and the exploration of new analytical capabilities. The chapter will close with concluding remarks on the transformative potential of these intelligent analytical systems for the future of the seafood industry.

\section{Achievement of Research Objectives}

The overarching goal of this thesis was to develop and validate a suite of machine learning methods to enhance the analytical capabilities of Rapid Evaporative Ionization Mass Spectrometry (REIMS) for marine biomass analysis, targeting rapid, automated, and \textit{in situ} application in fish processing plants.
This objective has been comprehensively achieved across three distinct research thrusts, successfully moving the field beyond traditional, limiting statistical methods (like OPLS-DA) toward a robust deep learning paradigm---a transition that is only just beginning in the wider food science literature \citep{Cardoso2025}.

\subsection*{1. Foundational Capabilities: Species and Body Part Identification (Chapter 4)}

The thesis successfully established a highly accurate, automated foundation for quality control. Novel Transformer and Mixture of Experts (MoE) Transformer architectures were developed and shown to consistently outperform OPLS-DA and all traditional machine learning methods. For the critical task of fish species classification, the MoE Transformer achieved a perfect 100\% accuracy, demonstrating that complex, sequence-aware deep learning models can fully exploit the chemical richness of REIMS data. Furthermore, by integrating Explainable AI (LIME and Grad-CAM), the research ensures that these high-performance models are not ``black boxes," providing chemically verifiable insights that are essential for industrial adoption.

\subsection*{2. Advanced Safety and Fraud Detection (Chapter 5)}

This research successfully formalized and solved novel contamination problems crucial for food safety. The thesis demonstrated that deep learning models, particularly those leveraging unsupervised pre-training (Masked Spectra Modeling) and Transfer Learning, are essential for handling subtle, difficult-to-detect signals:
\begin{itemize}
\item The models achieved competitive performance in the challenging ordinal multi-class oil contamination detection task (up to $49.10\%$ test accuracy), a significant improvement over OPLS-DA ($26.43\%$).
\item Crucially, this section established that the more challenging contamination tasks consistently benefited from positive knowledge transfer, providing practical guidance for accelerating model development in industrial settings.
\item By comparing ordinal classification techniques, the thesis validated that models which respect the inherent order of contamination levels (e.g., CORAL/CLM) yield more meaningful predictions by minimizing error distance.
\end{itemize}

\subsection*{3. Novel, Label-Free Batch Traceability (Chapter 6)}

The thesis achieved the critical objective of developing a practical, chemical-based solution for batch traceability, overcoming the limitations of costly physical tagging \citep{Mai2010} and complementing the field's current focus on vulnerable, data-entry-dependent digital traceability systems \citep{turkson2025digital, jiang2025traceability}. The novel SpectroSim framework, which utilizes a self-supervised contrastive learning approach with a Transformer encoder, achieved a remarkable $70.8\%$ balanced accuracy in determining if two samples shared a batch origin. This result:
\begin{itemize}
\item Vastly outperformed traditional binary classification methods (which struggled to surpass $60\%$ accuracy).
\item Eliminated the need for explicit class labels during training, making it a practical, plug-and-play solution for rapid, cost-effective deployment in processing environments, thereby addressing the high-cost and infrastructure barriers that hinder the adoption of digital-only systems \citep{Untal2025}.
\end{itemize}

\vspace{12pt} % Space before the summary paragraph

In summary, by systematically advancing the analytical capabilities of REIMS data, from fundamental species identification to complex trace contamination detection and label-free traceability, this thesis provides a validated suite of highly accurate, efficient, and potentially interpretable machine learning methods. The foundational goal of enabling rapid, automated, and \textit{in situ} quality assurance in the fish processing industry has therefore been successfully realized.

\section{Summary of Key Findings}

In this thesis, we have demonstrated that deep learning can be applied to REIMS marine biomass data to enhance the performance of ML analysis. The results have shown that deep learning improved classification and pair-wise comparison performance when compared to traditional ML methods. This section discusses the key conclusions of each research objective covered in chapters (\Cref{chapter:fish-species-and-part-identification,chapter:oil-contamination-and-cross-species-adulteration,chapter:contrastive-learning-for-batch-detection})

The core quantitative achievements of this thesis are distributed across Chapters 4, 5, and 6. To provide a consolidated overview, Table \ref{tab:results_summary} synthesizes the main experimental outcomes. This table presents the best-performing model developed for each of the five analytical tasks, directly comparing its balanced test accuracy against the traditional OPLS-DA benchmark to highlight the significant performance gains achieved by the proposed deep learning frameworks.

\begin{table*}[t]
\small 
\centering
\caption{Consolidated Summary of Best Model Performance Across All Tasks}
\label{tab:results_summary}
\begin{tabular}{l l l l l}
\toprule
\textbf{Analytical Task} & \textbf{OPLS-DA} & \textbf{Proposed Best Model} & \textbf{Best Model} & \textbf{Gain} \\
\midrule
Fish Species & 96.39\% & MoE Transformer & 100.00\% & +3.61\% \\
Fish Body Part & 51.17\% & Ensemble Transformer & 74.13\% & +22.96\% \\
Oil Contamination & 26.43\% & TL MoE Transformer & 49.10\% & +22.67\% \\
Cross-species Adulteration & 79.96\% & Pre-trained Transformer & 91.97\% & +12.01\% \\
Batch Detection & 53.19\% & SpectroSim (Transformer) & 70.80\% & +17.61\% \\
\bottomrule
\multicolumn{5}{l}{\footnotesize{All accuracy values are balanced test accuracy reported in their respective chapters.}}
\end{tabular}
\end{table*}

\subsection{On Foundational Classification (Chapter 4)}

In \Cref{chapter:fish-species-and-part-identification}, we proposed Transformers, both with and without unsupervised pretraining, for fish species and fish body part identification. Here, we draw the attention of the reader to two key findings of this chapter:
(1) deep learning supremacy over traditional ML models,
(2) the benefits of explainable AI techniques over black-box deep learning models,
and (3) balanced utilization for the MoE Transformer.
% and (4) positive transfer effects for more difficult tasks.

\subsubsection{Deep Learning Supremacy}

This thesis finds that across both tasks, fish species and fish body parts identification, the deep learning methods outperform traditional ML techniques and the incumbent OPLS-DA. This finding is not surprising; traditional models are simpler and often impose stronger inductive bias. A stronger inductive bias means the model makes more assumptions about the nature of the data. Simpler models (e.g., OPLS-DA, linear regression, and decision trees) have less capacity to learn complex patterns, even though OPLS-DA remains a robust benchmark for binary fish speciation \citep{Black2017, Shen2020, Shen2022}. Whereas more complex models, like those found in deep learning, have weaker inductive biases, they make fewer assumptions about the underlying nature of the data, enabling the capability to learn much more complex patterns. The marine biomass REIMS data is inherently structured, or sequential data, with long-range dependencies and complex feature interactions that simpler ML models are unable to grasp, an issue that has led to mixed results in the literature where conventional ML models like SVM or RF have sometimes failed to outperform traditional chemometrics \citep{De2023}. It takes those more complex models, capable of capturing complex patterns in structured data, to effectively harness the full potential of REIMS data.

\subsubsection{Explainable AI}

However, these more complex models with weaker inductive biases come at a cost. They are often referred to as black-box, as gradient-descent methods of models with millions of parameters of floating-point matrices, with no clear explanation or method to the madness, for why they make decisions. Deep neural networks (e.g., our proposed Transformer that consists of four encoder layers) capture complex feature interactions between intermediate layers that are safe to consider beyond human comprehension. To offer some respite, post-hoc (e.g., applied after the black-box model is trained) interpretability methods such as LIME and Grad-CAM impose strong inductive biases on the black-box method and make a series of simplifying assumptions to generate human-understandable explanations. Grad-CAM assumes that for the gradient information flowing into the final layer, features with larger gradients contributing to a particular class are more important to a decision. LIME approximates a black-box model's behaviour with a simple interpretable model (e.g., sparse linear model) within the local vicinity of the instance being explained. Some level of fidelity, or complexity, is lost in generating human-understandable explanations for black-box models, but these explanations help to bridge the gap between deep learning and human knowledge, addressing a key challenge identified in recent literature on applying ANNs to REIMS data \citep{Cardoso2025}.

\subsubsection{Balanced Expert Utilization}

Another key finding of this thesis was that the experts had balanced utilization across all four classification tasks. This indicates the effective use of model capacity without expert collapse. Expert collapse happens when only a small subset of available ``expert'' networks are consistently used or trained, while others are mostly ignored or redundant. Essentially, instead of all experts learning distinct aspects of the data, a few experts end up doing most or all of the work. Our results show the experts don't suffer from expert collapse, which is a sign the MoE is functioning as intended. This is how an MoE is supposed to work. Multiple expert networks and a gating network constitute an MoE model. The role of the gating network is to route different types of data to the most suitable expert(s). The idea is that each expert specializes in a particular subdomain of the input space, leading to a more efficient and powerful model. Our results demonstrate that the experts avoided collapse, a critical factor for an MoE functioning as intended. Crucially, this balanced utilization directly contributed to the high levels of accuracy and robustness observed across all four classification tasks.
This balanced utilization was pivotal because it allowed diverse experts to specialize in distinct aspects of the REIMS spectral signatures, directly enhancing the model's ability to differentiate accurately across all four classification tasks. Such effective use of the model's capacity, preventing any single expert from being overwhelmed or others from becoming redundant, underpins the strong predictive performance and reliability demonstrated in this work.

\subsection{On Contamination and Adulteration (Chapter 5)}

In \Cref{chapter:oil-contamination-and-cross-species-adulteration}, we extended the previous contribution to two new tasks of Oil Contamination and Cross-species Adulteration detection, finding that the proposed Transformer model with unsupervised pertaining significantly outperformed the traditional ML methods and other deep learning methods. In the following section, we highlight two key findings from this analysis:
(1) pretraining strategies work for mass spectra,
(2) Oil Contamination Detection is hard,
% (3) MoE Transformer supremacy,
and (3) Oil contamination consistently benefits from transfer learning.

\subsubsection{Pretraining Strategies}

This thesis finds that unsupervised pretraining strategies, such as the proposed Masked Spectra Modeling (MSM), offer consistent performance gains for oil contamination and cross-species adulteration detection. This suggests that there is inherent value in representation learning from mass spectra data in the task of BERT-styled left-to-right progressive masking. The task may seem simple, to predict which spectra come next, but since the data is structured and not tabular, and there is an inherent sequence to the features on the x-axis, in order of ascending mass-to-charge ratios, a sequence that can be learned. An intentional bonus side effect of this task is that the model implicitly learns to make peak predictions. In its journey from left to right, it will encounter several peaks for which it has to predict accurately. This implicitly learned task of peak prediction mirrors the manual analytical chemistry techniques that chemists practice in a lab when analyzing mass spectrometry samples by hand. This is analogous to the traditional analytical methods and is enhanced by its unsupervised nature; we no longer require a human-in-the-loop to manually annotate peaks. This saves a significant amount of time and domain expertise and enables the automation of REIMS marine biomass analysis.

\subsubsection{Oil Contamination Detection is Hard}

Another key finding of this thesis was that oil contamination detection is hard. Not only is it a multi-class classification task with seven classes. It can technically be framed as an ordinal classification task, where overlap exists between the various concentrations of oil contamination (e.g., 50\%, 25\%, 10\%, 5\%, 1\%, 0.1\%, 0\%), making the distinctions between individual classes more opaque. The inherent difficulty arises from the need to preserve and utilize this inherent order, which standard classification methods might ignore. Adjacent categories in an ordinal scale can often have ambiguous boundaries. Not only is this an ordinal classification problem, but it is also a multi-class classification task, with an increased number of decision boundaries that each model must formulate. With more classes, there is a higher chance of misclassification. The aforementioned ordinal scale of the oil concentrations adds inter-class similarity, making it more difficult for the model to learn fine-grained distinctions. The classification performance across all methods of oil contamination detection demonstrated the difficulty of this task.

\subsubsection{Oil Contamination Consistently Benefits}

This thesis finds that oil contamination detection consistently benefits from Transfer Learning. Our results show that regardless of the source domain, the target domain of the oil contamination dataset has positive transfer effects. Oil contamination gives a set of complex chemical signals that have underlying characteristics (e.g., they are primarily composed of hydrocarbons) that are fundamental enough to be recognized by models trained on diverse REIMS datasets of marine biomass. The knowledge gained from analyzing the general chemical landscape of marine biomass in a given source domain, e.g., fish species, fish body parts, or cross-species adulteration, provides a strong starting point for isolating the exogenous (i.e., foreign) signals of oil. The finding that positive transfer is observed regardless of the source domain suggests the learned features are capable of capturing fundamental aspects of mass spectrometry data from marine biomass, after which the specific oil signatures can easily be identified. Practically, this means that developing effective oil contamination detection models can be accelerated and improved by utilizing pre-existing models trained on other marine biomass datasets, potentially reducing the need for extensive de novo data collection and annotation for this specific, yet critical, application. Beyond oil contamination, these findings may encourage the exploration of transfer learning for other trace contaminant detection or anomaly identification tasks within REIMS analysis of complex biological matrices.

\subsection{On Batch Traceability (Chapter 6)}

In \Cref{chapter:contrastive-learning-for-batch-detection}, we demonstrated that Contrastive Learning, with substantial modifications, could be effectively applied to batch detection, exceeding the performance of traditional binary classification. Below, we highlight two key findings from this work: (1) deep learning supremacy, and (2) contrastive learning outperforms binary classification.

\subsubsection{Deep Learning Supremacy}

A key finding of this thesis was that deep learning outperformed traditional machine learning methods. Again, this can likely be put down to the structured sequential nature of the data, and this complexity not being captured by simpler models that impose a high inductive bias. Deep learning methods are simply better at handling complex sequential data. But let's explore further, deep learning methods have some key advantages revolving around their ability to automatically learn complex representations and handle intricate relationships in the data. Traditional machine learning methods rely heavily on manual feature engineering. If the chosen features (i.e., difference vector between pairs) are not discriminative enough, the model's performance will be limited. However, deep learning models, especially Siamese networks, can automatically learn hierarchical feature representations directly from the raw data. The model learns to identify the most relevant features for comparing items without explicit instructions (i.e., representation learning). Many traditional machine learning models (e.g., logistic regression, SVM) are inherently linear or assume specific data distributions. While some models (e.g., kernelized SVMS or random forests) can capture non-linearities, they are not as flexible as deep neural networks. Neural networks are universal approximations \citep{Hornik1989}, with multiple layers and non-linear activation functions. This means they can capture highly complex, non-linear relationships between paired items, a limitation recently demonstrated in food science where ANNs outperformed traditional PLS-DA by better modeling complex, non-linear relationships \citep{Cardoso2025}. This matters when the similarity is not additive or a linear function of the features, but involves complex feature interactions.

\subsubsection{Contrastive Learning Outperforms Binary Classification}

This thesis demonstrates that contrastive learning is superior to binary classification for batch detection of REIMS marine biomass. The difference vector approach, where we calculate the similarity of the pair as the difference vector between their features and train a standard binary classifier on these difference vectors to predict if the pair is similar or dissimilar. The difference vector is a crude distance measure that doesn't capture the full inherent semantic features of similarity more broadly. Potentially, the daw difference is not the most informative way to capture complex relationships in the high-dimensional REIMS data. In contrast, contrastive learning aims to generate an embedding space where similar items are closer together, and dissimilar items are further apart. The encoder network learns to produce these embeddings, or learned representations, of dissimilar and similar pairs. Two samples that belong to the same batch are a positive pair, and samples from different batches form a negative pair. The NT-XENT loss function then works by pulling positive pairs close together and pushing negative pairs apart. By contrasting many positive and negative examples, the model learns which features are invariant to positive pairs and which features distinguish between negative pairs.
The resulting embedding space facilitates nuanced similarity assessment, using the measure of cosine similarity, which effectively captures semantic relationships by considering the angle between feature vectors, irrespective of their magnitude. This provides a robust, intrinsic verification method, which is a critical component missing from modern, logistics-focused digital traceability systems \citep{turkson2025digital, dahariya2025enhancing}.

\section{Limitations of the Research}

Despite encouraging results, certain limitations of this research should be acknowledged:

\begin{itemize}
\item \textbf{Dataset Scope:} While the datasets from AgResearch provided a valuable foundation, they are specific to certain fish species (primarily Hoki and Mackerel) and processing conditions within New Zealand. The generalizability of the findings to a broader range of species, geographical origins, and processing environments requires further investigation. For example, geographical origin is a complex variable in its own right, with recent REIMS and ML studies focusing specifically on authenticating the provenance of a single species across different regions \citep{Lu2024}.
\item \textbf{Sample Size for Specific Tasks:} Some datasets, particularly for fish body part classification, had a limited number of samples, which can affect model robustness and the statistical significance of some findings, although techniques were employed to mitigate this.
\item \textbf{Dynamic Environmental Factors:} The current research primarily utilized static datasets. Real-world processing environments can have dynamic factors (e.g., seasonal variations \citep{OpenSeas2021} in fish biochemistry, changes in processing equipment) that might affect REIMS spectra and model performance over time.
\item \textbf{Computational Resources:} While efforts were made to develop efficient models, some advanced architectures like large Transformers or extensive MoE configurations can still be computationally intensive, potentially posing a barrier for facilities with limited resources.
\end{itemize}

\section{Future Work}

The engagement with REIMS marine biomass analysis detailed in this thesis has illuminated several key research directions. This section aims to present these topics, identified during our work, which we assert are promising avenues for future investigation likely to yield significant impact.

\subsection{Ordinal Classification for Oil Contamination}

The oil contamination task, defined by seven ordered concentration levels, presents a classic ordinal classification challenge where minimizing prediction error distance, measured by the Mean Absolute Error (MAE), is critical. While initial experiments confirmed that ordinal-aware methods like CORAL and Cumulative Link Models (CLM) significantly reduced the MAE compared to nominal classification, substantial performance gains are likely achievable by exploring other advanced ordinal-aware models referenced in the literature. Future research should systematically evaluate models from the three main families of ordinal classification. This includes Threshold-Based Models, specifically Support Vector Ordinal Regression (SVOR) \citep{herbrich1999support, chu2007support}, which has demonstrated top-tier performance on complex tabular datasets by using large-margin parallel hyperplanes. The investigation should also incorporate advanced Deep Learning Ordinal Losses, such as the Earth Mover's Distance (EMD) \citep{dematos2019emd} and Conditional Ordinal Regression for Neural Networks (CORN) \citep{shi2021corn}. These loss functions directly penalize misclassifications based on their rank distance, offering a principled way to optimize the network's learning objective for the specialized ordinal metrics of MAE and Quadratic Weighted Kappa (QWK).

Furthermore, a deeper analysis of Binary Decomposition Approaches should be conducted. Specifically, implementing techniques that leverage the cumulative link approach, as seen in Ordinal CNN architectures \citep{Niu2016}, could offer a robust alternative. This involves replacing the network's final layer with $k-1$ sigmoid neurons to predict cumulative probabilities, thereby structurally enforcing the ordinal constraint. Given that the sequential nature of REIMS data aligns well with Transformer-based models, these ordinal losses and architectures should be integrated into the existing Transformer framework. This dedicated focus on advanced ordinal methodologies is expected to yield substantial improvements in the classification accuracy of the oil contamination dataset, ensuring that the model's predictions are highly precise and chemically meaningful by minimizing large rank errors.

\subsection{Real-World Validation and Deployment}

The developed models should be deployed in their real-world application, a fish processing factory, where they are tested and validated in real-time fish processing to assess their capabilities.
The fish species and body part identification model developed in \Cref{chapter:fish-species-and-part-identification}, the oil contamination and cross-species adulteration detection model developed in \Cref{chapter:oil-contamination-and-cross-species-adulteration}, and the batch detection model developed in \Cref{chapter:contrastive-learning-for-batch-detection}, can all be deployed in a real-world testbed to measure their performance. Obviously, for real-world adoption by fish processing plants, the model would need to meet or exceed the current benchmarks for quality assurance and food safety. The future CyberMarine flex-factory \citep{PFR2020} serves as a perfect testbed to deploy and evaluate these models, since it already has the REIMS equipment, and processed the existing samples that formulate the dataset this thesis was based on, so it has access to the fish species, body parts, and contaminants used to replicate the experimental settings of the dataset.

For batch traceability, the current state-of-the-art approach that involves using RFID chips to track individual batches is cumbersome, impractical, and expensive \citep{Mai2010}.
After an initial investment in REIMS equipment, a relatively inexpensive, unintrusive, and plug-and-play solution to traceability is REIMS + SpectroSim --- the method for batch traceability through contrastive learning developed in \Cref{chapter:contrastive-learning-for-batch-detection}. This is a cheap, straightforward, and practical way to address the failings of traditional traceability in fish processing \citep{Helyar2014}. Furthermore, it serves as a critical analytical verification layer for the current state-of-the-art digital traceability systems \citep{Gastaldi2025}, which are often based on blockchain \citep{dahariya2025enhancing} or Digital Product Passports \citep{jiang2025traceability}. These digital systems are vulnerable to fraudulent data entry \citep{turkson2025digital}, a gap that ``SpectroSim" directly addresses. By providing a low-cost, label-free chemical check, it helps overcome the practical adoption barriers, such as cost and fisher-side infrastructure, that currently hinder digital-only solutions \citep{Untal2025}. Furthermore, once trained, the model does not require labels during inference to detect if fish belong to the same batch with 70.8\% accuracy. So the manual labor and the subsequent slowdown to supply chains, of adding physical labels to each batch, can be automated away entirely.

\subsection{Online Learning and Model Adaptation}

Investigating online learning or adaptive model retraining strategies would be crucial for maintaining model performance in dynamic industrial settings where data distributions might shift over time. One such example of a data distribution shift is the seasonal variation of Hoki. During spawning season, when Hoki are laying their eggs, the mother extracts all their lipids and puts them into their eggs as nutrients for them \citep{OpenSeas2021}. This dramatically changes the chemical composition of a mother Hoki, demonstrating a seasonal variation that would result in a data distribution shift for adult females of the Hoki population. We need machine learning models that are robust to such data distribution shifts, such as seasonal variation. Two approaches facilitate the change needed to handle data distribution shifts.

Online learning \citep{Hoi2021} is a machine learning approach where the model learns sequentially from incoming data instances, updating its parameters incrementally rather than retraining on the entire dataset from scratch. Model adaptation, often a goal or characteristic of online learning, refers to the process by which a pre-trained model adjusts its parameters to new data or a changing environment, allowing it to maintain performance or learn new patterns over time.
Online learning is a machine learning approach where the model learns incrementally by processing data sequentially, one instance or a small batch at a time, allowing it to continuously update and evolve. Model adaptation \citep{Chougiouris2024,Kabra2024} is the process by which a machine learning model adjusts its parameters or structure in response to new data or changes in the data distribution, often employing online learning techniques to maintain performance in dynamic environments.
Both of these approaches could be used to extend REIMS marine biomass analysis to handle data distribution shifts like seasonal variation.

\subsection{Multi-Modal Data Fusion}

Integrating REIMS data with other sensory information (e.g., nuclear magnetic resonance spectroscopy \citep{Bettjeman2018}, DNA barcoding \citep{Khaksar2015}) could create more comprehensive and resilient authentication systems. In machine learning, this process of combining data from multiple modalities into a single model is referred to as multi-modal data fusion \citep{Lahat2015}. In this case, different types of data sources are called modalities (e.g., REIMS, NMR, DNA barcoding). Multi-modal data fusion is the process of combining information from multiple, different modalities to create a more comprehensive, accurate, and robust understanding of a phenomenon or to improve the performance of a task \citep{Zhao2024}. Kaiser et al. \citep{Kaiser2017} demonstrate one such example of multi-modal data fusion, with their multi-modal transfer learning, that shares knowledge between the modalities of text, image, and audio. However, this is just one example of multi-modal data fusion; others include feature-level fusion, decision-level fusion, hybrid fusion, attention fusion, and so on.

Future work should consider the application of multi-modal fusion to REIMS marine biomass analysis to enhance performance.  Initial investigations could explore how different fusion strategies—such as early concatenation of REIMS spectral features with derived NMR chemical descriptors, or late fusion of probabilistic outputs from separate REIMS and DNA barcoding classifiers—impact specific tasks like fine-grained species differentiation or the detection of subtle adulteration. Furthermore, exploring hybrid approaches or attention mechanisms could be particularly fruitful for adaptively weighing the information from REIMS against complementary modalities like DNA barcoding, especially when dealing with samples of varying quality or complexity. This aligns with emerging research that combines different analytical targets, such as mercury levels, with machine learning to enhance food safety and traceability monitoring \citep{Piroutkova2025}.

\subsection{Advanced Model Architectures and Hybrid Approaches}

Continued exploration of novel deep learning architectures, including graph neural networks for modelling molecular relationships or hybrid models combining the strengths of different approaches (e.g., CNN-Transformer \citep{LeCun1989a,Vaswani2017} for spatiotemporal spectral analysis), could yield further improvements.
Graph Neural Networks (GNNs) \citep{Wu2020} are a class of neural networks designed to perform machine learning directly on graph-structured data, capturing relationships and dependencies between entities represented as nodes and edges. They work by iteratively aggregating information from a node's neighbors, allowing the network to learn representations (embeddings) of nodes, edges, or entire graphs for tasks like node classification, link prediction, or graph classification.

Future REIMS marine biomass analysis could greatly benefit from Graph Neural Networks (GNNs). These can model spectral features (m/z values) as graphs—with edges indicating chemical relationships (e.g., neutral losses, isotopic patterns) or learned correlations—to exploit inter-feature dependencies for more robust sample classification or biomarker discovery. Alternatively, GNNs can represent individual samples as nodes in a graph, with edges denoting similarities or spatio-temporal links, to learn holistic patterns across entire datasets. Exploring optimal graph construction methodologies and suitable GNN architectures (like Graph Convolutional or Attention Networks) for such analytical tasks is a key research direction. Furthermore, hybrid models promise enhanced capabilities by synergistically integrating architectures; for instance, combining CNNs (for extracting potent features from raw REIMS spectra) with GNNs (to subsequently model inter-sample relationships using these features or incorporate chemical knowledge graphs), or pairing Recurrent Neural Networks (RNNs/LSTMs) for analyzing time-series spectral data with GNNs to track evolving compound relationships and derive deeper insights.

\section{Final Remarks}

This thesis has demonstrated the profound impact that advanced machine learning, when thoughtfully applied to high-resolution REIMS data, can have on addressing critical challenges in the seafood industry. By systematically developing and evaluating a range of techniques from sophisticated transformer architectures to strategic transfer learning and innovative self-supervised methods, this research has not only pushed the boundaries of analytical capability but also provided practical pathways towards more secure, transparent, and efficient seafood supply chains. The journey from identifying a fish out of water to identifying a fish out of data has revealed that the fusion of cutting-edge spectrometry and intelligent algorithms is key to safeguarding one of our most valuable global food resources. The insights gained and the methodologies developed herein lay a strong foundation for future innovations in food science and analytical chemistry \citep{Xue2025}, promising a future where food fraud is more readily detected and food safety is more robustly assured.